\section{Generalizations to prime powers}\label{sec:generalization}

Catalan numbers may be generalized to $q$-Catalan numbers, which have the formula 
\[C^{(q)}_n=\frac{1}{(q-1)n+1}\binom{qn}{n}.\]
In this section, we will use $C^{(q)}_n$ to denote $q$-Catalan numbers and $C^{(q)}_n(b)$ to denote weighted $q$-Catalan numbers.

These count the number of paths from $(0, 0)$ to $(qn, 0)$ with steps $(1, q-1)$ and $(1, -1)$ which never go below the $x$-axis and the number of $q$-ary trees on $n$ vertices (each vertex has $q$ distinguishable branches). Similarly, there are the weighted $q$-Catalan numbers. For a $q$-ary tree $\mathcal{T}$, weight each vertex $v$ by $b(i)$ where $i$ is the number of non-right edges on the path from the root to $v$, and let the weight of $\mathcal{T}$ be the product of the weights of its vertices. The weighted $q$-Catalan numbers are the sums of the weights of all $q$-ary trees on $n$ vertices:
\[C_n^{(q)}(b)=\sum_{\mathcal{T}}\wt(\mathcal{T}).\]
The weighted $q$-Catalan numbers have generating function
\[\sum_{n\geq0}C_n^{(q)}(b)x^n=\cfrac{1}{1-\cfrac{b(0)x}{\left(1-\cfrac{b(1)x}{1-\left(\cfrac{b(2)x}{(1-\cdots)^{q-1}}\right)^{q-1}}\right)^{q-1}}}.\]
As before, there is a symmetry group $G_n$ on $q$-ary trees of depth $n$, generated by permuting the subtrees at any vertex. The complete $q$-ary trees are fixed under all symmetries and consist of layers of $1, q, q^2, \dotsc, q^{k-1}$ vertices.

Konvalinka \cite{konvalinka2007divisibility} generalized Postnikov and Sagan's result to divisibility of $q$-Catalan numbers when $q$ is a power of a prime. Similarly, we may extend the techniques for the case $q=2$ to prove the following result:

\begin{theorem}\label{thm:qmain}
Let $q=p^k$ be a prime power. If $b: \Z_{\geq 0}\to\Z$ satisfies 
\begin{enumerate}
    \item $b(0)\equiv 1\pmod{q}$
    \item $q^2\mid (\Delta b)(x)$ for all $x\in\mathbb{Z}_{\geq 0}$,
    \item $q^n\mid (\Delta^nb)(x)$ for all $n\geq 2$ and $x\in \mathbb{Z}_{\geq 0}$,
\end{enumerate}
then
$C_n^{(q)}(b))\equiv C_n^{(q)}\pmod{p^{\xi+k}}$
where
\[\xi = \frac{s_p((q-1)n+1)-1}{p-1}=\xi_p(C_n^{(q)}).\]
\end{theorem}

The proof is similar so we will show the main claims and only give brief details of the proofs. Define $\mathcal{F}$ as the set of functions satisfying $q^n\mid \Delta^n f$ for all $n$, and the map $f\mapsto (\varepsilon_0^f, \varepsilon_1^f, \dotsc)\in (\mathbb{Z}/q\mathbb{Z})^{\mathbb{N}}$ so that $\Delta^n f\equiv \varepsilon^f_nq^n\pmod{q^{n+1}}$.

\begin{lemma}
The set $\mathcal{F}$ is closed under the following operations:
\begin{align*}
    &f\mapsto Sf, \quad (f, g)\mapsto f\cdot g, \quad \text{ and } \\
    &(f_1,\dotsc, f_q)\mapsto \langle f_1, \dotsc, f_q\rangle = \frac{1}{q}\sum_{i=1}^{q}f_i(x+1)\prod_{j\neq i}f_j(x).
\end{align*}
\end{lemma}
\begin{proof}
    Closure under $S$ is clear, and closure under products follows from Lemma \ref{lem:productrule} (the product rule). Closure under $\langle \dotsb\rangle$ follows by writing
    \[\langle f_1, \dotsc, f_q\rangle =\frac{1}{q}\sum_{i=1}^{q}(f_i+\Delta f_i)\prod_{j\neq i}f_j,\]
    expanding, and using the product rule.
\end{proof}
We define average weight functions $r_b(\mathcal{O}; x)$ for orbits as before, and these weight functions are in $\mathcal{F}$. Now we give the analog of Lemma \ref{lem:epsilon formula}.

\begin{lemma}
\label{lem:qepsilonformula}
Let $\mathcal{O}$ be an orbit of $q$-ary trees and let the orbits of the subtrees of the root vertex be $\O_1, \dotsc, \O_q$. Then
\[\varepsilon^{\mathcal{O}}_{m} \equiv\sum_{i_1+\dotsb+i_q+k=m}\binom{m}{i_1, \dotsc, i_q, k}\varepsilon_{k}(\varepsilon^{\mathcal{O}_1}_{i_1}\dotsm\varepsilon^{\mathcal{O}_q}_{i_q}+\varepsilon^{\mathcal{O}_1}_{i_1+1}\dotsm\varepsilon^{\mathcal{O}_q}_{i_q}+\dotsb+\varepsilon^{\mathcal{O}_1}_{i_1}\dotsm\varepsilon^{\mathcal{O}_q}_{i_q+1})\pmod{q}.\]
(The expression inside the parentheses contains $q+1$ terms, one term with indices $i_1, \dotsc, i_q$ while the other $q$ have exactly one index incremented.)
\end{lemma}
\begin{proof}
We use the recursive formula
\[r_b(\mathcal{O})=b\cdot \langle r_b(\O_1), \dotsc, r_b(\O_q)\rangle.\]
By expanding $\langle f_1, \dotsc, f_q\rangle =\frac{1}{q}\sum_{i=1}^{q}(f_i+\Delta f_i)\prod_{j\neq i}f_j$, we arrive at the given formula.
\end{proof}
Before stating the analog of Lemma \ref{lem:reduction}, we first have the following proposition:

\begin{proposition}
\label{prop:binomialdivisibility}
Let $q$ be a positive integer and $i_1, \dotsc, i_q$ be nonnegative integers, not all zero. Suppose that the distinct elements of the set $\{i_1, \dotsc, i_q\}$ are $j_1, \dotsc, j_t$ and these elements appear with multiplicity $k_1, \dotsc, k_t$. Then $q$ divides
\[\binom{i_1+\dotsc +i_q}{i_1, \dotsc, i_q}\binom{q}{k_1, \dotsc, k_t}.\]
\end{proposition}
\begin{proof}
    The binomial expression given counts the number of ways to assign $i_1+\dotsb+i_q$ distinguishable people into $q$ distinguishable rooms so that the distribution of the number of people in the rooms is a permutation of $i_1, \dotsc, i_q$. Since there is at least one person, we can group valid assignments with cyclic shifts (move each person down one room), so the number of assignments it a multiple of $q$.
\end{proof}

\begin{lemma}
\label{lem:qreduction}
    Let $\mathcal{O}$ be an orbit of $q$-trees on $n$ vertices, and suppose that for some vertex $v$, the subtree with root $v$ is a complete $q$-ary tree of depth $k$. Let $\mathcal{O}'$ be the same orbit of trees with the subtree at $v$ replaced by a single vertex at $v$. Then \[\varepsilon_m^{\mathcal{O}}=\varepsilon_0^{\frac{q^k-1}{q-1}-1}\varepsilon_m^{\mathcal{O}'}.\]
\end{lemma}
\begin{proof}
    Like the proof of Lemma \ref{lem:reduction}, it suffices to analyze the case where $\O$ is a complete $q$-ary tree and $v$ is its root. We expand with Lemma \ref{lem:qepsilonformula}. Since the orbits $\mathcal{O}_1, \dotsc, \mathcal{O}_q$ are the same, we have
    \[\sum_{i_1+\dotsb+i_q+k=m}\binom{m}{i_1, \dotsc, i_q, k}\varepsilon_{k}(\varepsilon^{\mathcal{O}_1}_{i_1+1}\dotsm\varepsilon^{\mathcal{O}_q}_{i_q}+\dotsb+\varepsilon^{\mathcal{O}_1}_{i_1}\dotsm\varepsilon^{\mathcal{O}_q}_{i_q+1})=0\pmod{q}.\]
    The remaining term is\[\varepsilon^{\mathcal{O}}_{m} \equiv\sum_{i_1+\dotsb+i_q+k=m}\binom{m}{i_1, \dotsc, i_q, k}\varepsilon_{k}\varepsilon^{\mathcal{O}_1}_{i_1}\dotsm\varepsilon^{\mathcal{O}_1}_{i_q}\pmod{q}.\]
    (We have replaced all $\O_i$ with $\O_1$.) Given a multiset of indices $\{i_1, \dotsc, i_q\}$, with distinct elements $j_1, \dotsc, j_t$ that appear $k_1, \dotsc, k_t$ times, the number of times the term $\varepsilon^{\mathcal{O}_1}_{i_1}\dotsm \varepsilon^{\O_1}_{i_q}$ appears in the above sum is $\binom{q}{k_1,\dotsc, k_t}$, and it is multiplied by $\binom{m}{i_1, \dotsc, i_q, k}$. Since $\binom{i_1+\dotsb +i_q}{i_1, \dotsc, i_q}$ divides $\binom{m}{i_1, \dotsc, i_q, k}$, by Proposition \ref{prop:binomialdivisibility}, these terms cancel out unless $i_1=\dotsb=i_q=0$. Hence we are left with
    \[\varepsilon^{\O}_m=\varepsilon_m(\varepsilon^{\O_1}_0)^q\]
    and the statement of the lemma follows from there.
\end{proof}

The proof of Theorem \ref{thm:qmain} is quite similar to the proof of Theorem \ref{thm:main} from here. We construct all minimal orbits on $n$ vertices according to the analog of Theorem \ref{thm:orbit structure}, and each one may be reduced using Lemma \ref{lem:qreduction} to get orbits of trees on a smaller number of vertices. This will reduce to the case of trees on a smaller number of vertices until we have reduced to $n\leq q$. In these cases, we directly verify that $\varepsilon_0=1$ and $\varepsilon_1=0$ is sufficient to have the sum of $\varepsilon^{\O}_0$ equal to $1\pmod{q}$.